\section{Bối cảnh và nền tảng đề tài}
Trong các chuyến du thuyền, hoạt động vận hành thường bao gồm nhiều khâu liên quan đến lịch trình, dịch vụ, thanh toán và tương tác với hành khách. Mỗi chuyến đi có thể kéo dài nhiều ngày, đồng thời đòi hỏi sự phối hợp chặt chẽ giữa các bộ phận điều phối, vận hành và quản lý.

\begin{itemize}
    \item Trước khi khởi hành, ban tổ chức cần lên kế hoạch về lịch trình, cảng ghé thăm, hoạt động giải trí, năng lực phục vụ, phí tham gia và chính sách hủy dịch vụ.
    \item Trong suốt chuyến đi, hành khách cần một phương thức thuận tiện để xem lịch trình, đăng ký hoạt động và phân biệt rõ ràng giữa dịch vụ miễn phí và trả phí.
    \item Do đặc thù kết nối mạng không ổn định trên tàu, hệ thống cần hỗ trợ ghi nhận giao dịch ngoại tuyến qua thiết bị POS, mã QR hoặc vòng tay RFID/NFC, và tự động đồng bộ dữ liệu khi có kết nối trở lại.
    \item Sau chuyến đi, bộ phận quản lý cần các công cụ mạnh mẽ để đối soát giao dịch, thanh toán hóa đơn cuối cùng, lập báo cáo vận hành và thu thập phản hồi khách hàng.
\end{itemize}